%%=============================================================================
%% CAPITULO 2 -- REFERENCIAL TEORICO
%% Esqueleto provisorio (dissertacao dummy): texto a ser desenvolvido.
%%=============================================================================

\section{Referencial teórico}

A dificuldade de descoberta e seleção de recursos em acervos volumosos é
problema recorrente na literatura de recomendação educacional, descrito como
sobrecarga de informação. \citeonline{castro2014} partem exatamente desse
problema ao propor a recomendação de objetos de aprendizagem para professores
que elaboram planos de aula: diante de repositórios extensos e de ferramentas
de busca pouco flexíveis, o docente tem dificuldade de localizar o material
adequado, o que torna a recomendação automatizada uma alternativa à busca
manual. É a mesma situação do professor da rede piauiense diante das videoaulas
do Canal Educação.

Os acervos educacionais públicos, por sua vez, aproximam-se da noção de
Recursos Educacionais Abertos, e a pesquisa sobre sua recomendação reconhece
que esses sistemas receberam atenção menor do que os voltados ao comércio
eletrônico e ao entretenimento, apesar de seu potencial para ampliar ganhos de
aprendizagem:

\begin{citacao}
Educational recommenders have received much less attention in comparison with
e-commerce- and entertainment-related recommenders, even though efficient
intelligent tutors could have potential to improve learning gains and enable
advances in education that are essential to achieving the world's
sustainability agenda. Through this work, we make foundational advances towards
building a state-aware, integrative educational recommender.
\parencite[p.~1]{bulathwela2022}
\end{citacao}

Esse desequilíbrio importa para a presente pesquisa por duas razões. Indica que
o problema de organizar e recomendar recursos educacionais permanece menos
desenvolvido do que problemas análogos em outros domínios, e aponta o início
frio, isto é, a escassez de dados nas etapas iniciais de uso, como desafio
agudo da recomendação educacional \parencite{bulathwela2022}. Como o artefato
proposto não dispõe de histórico de interação de docentes com o acervo,
abordagens baseadas em conteúdo e em similaridade semântica entre videoaula e
habilidade da matriz tendem a ser mais adequadas do que a filtragem
colaborativa.

\subsection{Representação semântica de videoaulas}

A construção de uma camada de articulação entre conteúdo e currículo depende de
representar semanticamente os recursos do acervo. A literatura demonstra que é
possível extrair automaticamente descritores de vídeos de aprendizagem e
utilizá-los para recomendar recursos correlatos, com resultados que se
aproximam das palavras-chave atribuídas por usuários humanos
\parencite{schulten2020}. Esse achado sustenta a viabilidade de mapear o
conteúdo audiovisual do Canal Educação, por meio da transcrição automática do
áudio, às habilidades da matriz de referência do ENEM.

As três obras convergem para a lacuna que o artefato enfrenta: nenhuma delas
toma um descritor curricular oficial como alvo da classificação nem confronta o
resultado com referência externa publicada por órgão avaliador. A pesquisa
adota, por isso, a atribuição oficial de habilidades a itens do exame,
divulgada pelo INEP nos microdados do ENEM, como referência independente de
quem construiu o classificador.
